\subsection{Type 12. Two cubic-field curves of rank one}\label{detail:type12}

\subsubsection{The two parity classes}\label{arith:61}

Write \(b\) for the square of the length-four B center.
A valid type-12 lift gives a point \(mP_1+nP_9\) with \(m+n\) odd on the
rank-two elliptic curve
\[E_{12}:Y^2=X^3-240X^2+215aX+90a^2,\qquad a=160,\quad X=ab,\]
where \(P_1=(160,-2400)\) and \(P_9=(1440,50400)\).

\emph{Proof.} From \S\ref{app:type12}, \(b\) and \(10-b\) are squares. The three
singleton squares have product \(f(b)/64\), where
\(f(b)=160b^3-240b^2+215b+90\). Indeed their monic cubic, multiplied
by 64, is
\[64Z^3-(240b+320)Z^2+(320b^2+280b+436)Z-f(b).\]
Thus \(f(b)\) is a square, giving the point with \(X=ab\).
The full Mordell--Weil group is free on \(P_1,P_9\).

Let \(w^3+38w-504=0\) and \(r=(4-w)/8\), so \(f(r)=0\).
At the ramified prime of \(\mathbf Q(w)\) over two, take the integer
valuation with \(v(2)=2\). We have \(v(w)=1,\ v(r)=-5,\ v(a)=10\).
The equation \(b+(10-b)=10\) expresses 10 as a sum of two rational
squares. Their roots must both be dyadic units. Indeed unequal
valuations give an even valuation for the sum, while two roots of equal
valuation \(j\) give valuation \(2j+1\), since two odd squares sum to
2 modulo 8. Hence \(v(b)=0\). The Kummer character at the root \(ar\)
is represented by \(a(b-r)\), whose valuation is five. The same holds
at both generators (with \(b=1,9\)). Valuation modulo two is additive
on the Mordell--Weil group, so \(m+n\equiv1\pmod2\).

For a surviving class indexed by \(b_0=1\) or \(9\), put \(d=b_0-r\).
The same abscissa character makes \((b-r)/d\) square. Multiplying by the
two center squares gives a point with \(X=-db\) on
\[Y^2=X(X+10d)(X+rd).\]
The two-torsion has order four, and unconditional two-descent gives
rank upper bounds three and four for the two covers, respectively.
\cert{check_type12_center_parity.m}

\subsubsection{Two rank-one curves and their rational images}\label{arith:62}

Suppose that \(b,10-b\) are rational squares and the
singleton cubic in \S\ref{arith:61} splits into three rational squares,
counted with multiplicity. Then \(b\in\{1,9\}\). Both values give a
collision. Type 12 has no valid lift at any rational height.

\emph{Proof.} Retain \(K,r,f\) from \S\ref{arith:61} and write
\[q(B)=f(B)/(160(B-r))=B^2+q_1B+q_0,\quad
 q_1=r-3/2,\qquad q_0=r^2-3r/2+43/32.\]
The proof of \S\ref{arith:61} uses just the square conditions in the
statement. Thus for one of \(b_0=1,9\), the quantity \((b-r)/(b_0-r)\)
is a square in \(K\). Since \(f(b)\) is square, \(160(b_0-r)q(b)\)
is also square. The rational number \(b\) is nonzero since 10 is
not a rational square.

For \(b_0=1\), put \(d=-160(1-r)\). The product \(d(b-10)q(b)\) is a
square. With \(x=d(b-10)\) it gives a point on
\[E^{(1)}:\quad y^2=x(x^2+dq'(10)x+d^2q(10)).\]
For \(b_0=9\), put \(D=160(9-r)\). This time take the square
\(D b(10-b)q(b)\). Division by \(b^4\) gives the square
\(D(10t-1)(q_0t^2+q_1t+1)\) with \(t=1/b\). Put \(c=10Dq_0\) and
\[F(U)=(U-c/10)(U^2+10Dq_1U+10Dc).\]
Then \(F(ct)=c^2D(10t-1)(q_0t^2+q_1t+1)\). We obtain a point on
\(E^{(9)}:y^2=F(x+c/10)\) with \(x=c/b-c/10\). In each case let
\(P_0\in E^{(b_0)}(K)\) be a point over \(b=b_0\).

\textbf{Ranks.} For the first curve do two-descent directly. For the second
do it on the quotient by the order-two subgroup generated by \((0,0)\)
and use the image of \(P_0\). Each descent gives a 2-Selmer group of
order 16 and rational two-torsion of order 2. The quotient of the
Selmer group by torsion and the specified point has dimension two.
On this quotient the alternating Cassels--Tate pairing takes the
nonzero value on a pair of basis elements. The image of the global
points lies in the radical, so the rank is at most one. Each \(P_0\)
has infinite order, and isogeny preserves rank. Thus both original
curves have rank one over \(K\).

The second descent uses the Minkowski proof bound throughout. For the
first, the quadratic cofactor of the two-division algebra has absolute
field \(L\) defined by \(T^6+228T^4+12996T^2+7077920\). PARI verifies
its discriminant \(-5540347111683392000\), checks that every ideal of
its factor base has norm at most 100, and certifies that the factor base
generates the class group. Magma uses the resulting proof bound 100
after testing field isomorphism, and the Minkowski bound elsewhere.
\cert{check_type12_fields.gp}; \cert{check_type12_exclusion.m}; \href{https://pari.math.u-bordeaux.fr/dochtml/html-stable/General_number_fields.html\#bnfcertify}{PARI class-group certification}

\textbf{The residue disks.} Take \(p=907\). Since \(-1\) is a nonsquare,
a pole of \(b\) is impossible. Indeed \((10-b)/b\) is square and would
reduce to \(-1\). The field \(K\) splits completely at \(p\). For each
residue \(b\), test the two center squares, the three squareclasses
\((b-r_j)(b_0-r_j)\) at the three roots \(r_j\), and whether the
singleton cubic has three roots counted with multiplicity, each a square in
\(\mathbf F_p\). These reductions are necessary since the cubic
has unit leading coefficient. In each class the exact enumeration of
the 907 residues leaves only \(b\equiv b_0\pmod p\).

For the three embeddings \(j:K\hookrightarrow\mathbf Q_p\), let \(n_j\)
be the order of the reduced elliptic curve and write
\(\ell_j(R)=\log_E([n_j]R_j)\), using the differential \(dx/(2y)\).
The displayed models have good reduction, and \(y(P_{0,j})\) is a unit
at all three places. With \(x(b)=d(b-10)\) or \(c/b-c/10\), compute
\[L_j=\ell_j(P_0)/p\bmod p,\qquad
 V_j=n_j x'(b_0)/(2y(P_{0,j}))\bmod p .\]
For \(L_j\) the formal parameter \(-x/y\) of \([n_j]P_{0,j}\) suffices
modulo \(p^2\). The certificate gives the following rows in a common
ordering of the embeddings.

{
\begin{longtable}[]{@{}
  >{\raggedright\arraybackslash}p{(\linewidth - 6\tabcolsep) * \real{0.2500}}
  >{\raggedright\arraybackslash}p{(\linewidth - 6\tabcolsep) * \real{0.2500}}
  >{\raggedright\arraybackslash}p{(\linewidth - 6\tabcolsep) * \real{0.2500}}
  >{\raggedright\arraybackslash}p{(\linewidth - 6\tabcolsep) * \real{0.2500}}@{}}
\toprule\noalign{}
\begin{minipage}[b]{\linewidth}\raggedright
\(b_0\)
\end{minipage} & \begin{minipage}[b]{\linewidth}\raggedright
\((L_1,L_2,L_3)\)
\end{minipage} & \begin{minipage}[b]{\linewidth}\raggedright
\((V_1,V_2,V_3)\)
\end{minipage} & \begin{minipage}[b]{\linewidth}\raggedright
\(L_iV_j-L_jV_i\) for \(ij=12,13,23\)
\end{minipage} \\
\midrule\noalign{}
\endhead
\bottomrule\noalign{}
\endlastfoot
1 & 823, 146, 636 & 453, 716, 192 & 698, 516, 760 \\
9 & 236, 627, 838 & 800, 274, 656 & 238, 499, 300 \\
\end{longtable}
}

For any global point \(R\), rank one gives a rational \(k\) such that
\(\ell_j(R)=k\ell_j(P_0)\) at every place. This follows by clearing
denominators in the group modulo torsion. At each place a candidate
reduces to \(\epsilon_jP_{0,j}\) with \(\epsilon_j=\pm1\). Choose two
places \(i,j\) with the same sign. Set \(A_i=\ell_i(P_0)/p\) and likewise
for \(j\). The function \(A_i\ell_j(R(b))-A_j\ell_i(R(b))\) vanishes
at the anchor with those signs, and at any global candidate on these
two local branches. Its derivative at the anchor reduces to
\(\epsilon_i(L_iV_j-L_jV_i)\), a unit in every case.

To see uniqueness on the entire disk, put \(h=b-b_0\). The derivatives
of the logarithms are \(n_j x'(b)/(2y_j(b))\). As power series in \(h\)
they have integral coefficients, since the ordinates are units and
Hensel lifting takes their square roots with unit derivative. For the
reciprocal formula \(b_0=9\) is a unit as well. Integration gives terms
of degree \(k\) in \(h\) with coefficients in \(k^{-1}\mathbf Z_p\).
For \(h\in p\mathbf Z_p\) and \(k\ge2\),
\((k-1)v_p(h)>v_p(k)\) when \(h\ne0\). The linear term therefore has
strictly smaller valuation than every higher term. This proves
\(b=b_0\) in both classes.

The singleton squares at the two anchors are respectively
\(\{1,9,25\}/4\) and \(\{25,49,81\}/4\).
At \(b=1\) the B length-four chain centered at 1 contains \(1/2\); at
\(b=9\) the one centered at 3 contains \(3/2\). Both vertices belong to
the centered A chain. The exceptional branches \(b=0,1,4\) of the fourth
moment already give collisions by \S\ref{app:type12}. This proves the type exclusion.

\subsection{Type 1. Split descent and a further division}\label{detail:type1}

\subsubsection{The parity of the squared center}\label{arith:63}

Put \(b=v^2\) for the square of the length-two B
center in type 1, and set \(a=16384\) and
\[f_1(B)=aB^3+38672B^2+33144B+11025,\qquad
 E_1:Y^2=X^3+38672X^2+33144aX+11025a^2.\]
For any valid lift, the point with \(X=ab\) lies in one of two classes
of \(E_1(\mathbf Q)/2E_1(\mathbf Q)\), represented by the points above
\(b=0,1\). Fourteen of the sixteen global classes are excluded.

\emph{Proof.} By \S\ref{app:type1}, \(5b+4\) is square. The constant coefficient of the
singleton cubic of \S\ref{app:type1} is \(-f_1(b)\), after scaling its leading
coefficient to 64. Hence \(f_1(b)\) is square, and
\((ab,a\sqrt{f_1(b)})\) gives the claimed point. Changing its ordinate
sign has no effect modulo twice the group. The full group is torsion-free
of rank four. The certificate computes a basis and represents its sixteen
parity classes by sums with coefficients zero or one.

At each prime in the table below, reduce into
\(E_1(\mathbf F_p)/2E_1(\mathbf F_p)\). For integral \(b\), retain the
residues where \(b,5b+4\) are squares and all three singleton roots
are squares, allowing multiplicity. For a pole, \((5b+4)/b\) reduces
to 5 and is square. The three singleton squares divided by \(b\) are
squares and integral at \(p\), since their cubic after division by
\(b^3\) has integral coefficients and unit leading coefficient. They
must reduce to the roots of \(64z^3-1152z^2+7232z-16384\). Retain the
group identity for a pole whenever both necessary tests pass.

{
\begin{longtable}[]{@{}llllll@{}}
\toprule\noalign{}
Prime & 7 & 67 & 73 & 107 & 257 \\
\midrule\noalign{}
\endhead
\bottomrule\noalign{}
\endlastfoot
Global classes left after intersection & 8 & 6 & 5 & 4 & 2 \\
\end{longtable}
}

Exact rational division by two identifies these last classes with the
two stated representatives. For a root \(r\) of \(f_1\) the abscissa character
at \(ar\) now makes \((b-r)(b_0-r)\) square in \(\mathbf Q(r)\).
At 73 the polynomial \(f_1\) splits with distinct roots, and 5 is a
nonsquare, so there is no pole. Testing the character at all three roots
and imposing the center and singleton squares gives just \(b=b_0\)
modulo 73 in each class. Sections~\ref{arith:64} and~\ref{arith:65} determine the eligible
centers in these disks.
\cert{check_type1_parity.m}

\subsubsection{Split descent in the cubic field}\label{arith:64}

In the type-1 parity class represented by \(b_0=1\) in
\S\ref{arith:63}, the only squared center satisfying the necessary square
and splitting conditions is \(b=1\). It gives overlapping chains. The
cubic-field elliptic cover for \(b_0=1\) has rank exactly one. For
\(b_0=0\) the corresponding cover has unconditional rank at most two.

\textbf{The two covers.} Put \(K=\mathbf Q(w)\), where
\(w^3-w^2-330w-149578=0\), and set
\(r=(-w^2+20w-5426)/7168\). Then \(f_1(r)=0\). By \S\ref{arith:63},
\((b-r)(b_0-r)\) is square in \(K\). Since \(b\) and \(5b+4\) are
rational squares, \(d b(b+4/5)(b-r)\) is square for \(d=5(b_0-r)\).
Consequently a lift gives a point on the following elliptic curve:
\[E_d:Y^2=X(X-dr)(X+4d/5),\qquad X=db .\]

\textbf{The global 2-Selmer groups.} Use the roots
\((e_1,e_2,e_3)=(0,dr,-4d/5)\). The split Kummer map has coordinates
\(X-e_i\) modulo squares, with their product square. At a two-torsion
point replace its zero coordinate by the product of the other two.
Take \(S\) to consist of all prime ideals over
\(2,3,5,7,11,12101,1122137\). Outside \(S\) the roots are integral with
unit differences and the residue characteristic is odd. There the
local image consists of the unramified triples of square product.
Indeed the image is unramified by good reduction, and both spaces have
dimension two. Thus the global classes lie in two copies of
\(K(S,2)\), the squareclasses of even valuation outside \(S\).

PARI certifies the class group and units of \(K\). The \(S\)-class
group is trivial, so the fundamental \(S\)-units and the torsion unit
span \(K(S,2)\). They give dimension 17. At each finite place in \(S\),
the certificate constructs independent Kummer images of local points
until their number is \(2+[K_v:\mathbf Q_2]\) at a dyadic place or
two at an odd place. These are the exact dimensions of
\(E_d(K_v)/2E_d(K_v)\), by the local Lie group structure and the four
two-torsion points. Local unit characters at \(v\mid2\) are computed
modulo \(\mathfrak p^{2v(2)+1}\), which suffices by Hensel's lemma.
At the real embedding the two sign patterns are given by the unbounded
component and by an abscissa between the two smallest roots.
Intersecting these linear conditions gives Selmer dimensions 8 and 7
for \(b_0=0,1\), respectively. The point selection stops only after
the required local dimension is attained.

\textbf{The Cassels--Tate pairing.} For a Selmer triple \(s=(s_1,s_2,s_3)\)
use the two-cover \(s_i z_i^2=X-e_i\). For cyclic indices \(i,j,k\)
let \(L_i\) be the tangent at a specified \(K\)-point of the conic
\[s_j z_j^2-s_k z_k^2+(e_j-e_k)t^2=0 .\]
Cassels' formula for the pairing against another triple \(s'\) is the
product over the places of the three Hilbert symbols
\((L_i(P_v),s'_i)_v\), where \(P_v\) is a local point of the cover.
The data file supplies four triples on each curve and three conic
points per triple. The verifier checks the conic identities and the
square-product conditions. It constructs local points at the bad
places, at primes supporting the triples and at primes supporting the
fractional ideal generated by the coefficients of each tangent.
At the remaining odd primes the reductions of the conics are smooth,
their tangent forms are primitive, and all entries of the triples
are units. On the conic a tangent cuts out twice its point of contact.
Thus the evaluation of a tangent on primitive integral coordinates
has even valuation whenever it is nonzero. The three conic coordinates
are primitive for a primitive point of the intersection, by the other
quadric equation. Changing projective scale has trivial total symbol
because the target triple has square product. This accounts for every
place. The verifier uses Hensel error bounds for the finite symbols
and rational isolating intervals for the real signs.
\href{https://www.dpmms.cam.ac.uk/~taf1000/papers/casselspairing.pdf}{Cassels pairing over number fields}

Using 1 for a negative symbol product, the resulting pairing minors
in the order of the data file are
\[\begin{pmatrix}0&1&0&0\\1&0&0&1\\0&0&0&1\\0&1&1&0\end{pmatrix},
 \qquad
 \begin{pmatrix}0&1&1&1\\1&0&0&1\\1&0&0&0\\1&1&0&0\end{pmatrix}.\]
Both are invertible over \(\mathbf F_2\). The images of global points
are in the radical. Subtracting four and then the two torsion
dimensions gives rank bounds two and one. For \(b_0=1\),
\(P=(d,3d^2/5)\) has infinite order, proving equality.
\cert{type1_descent/verify.gp},
\cert{type1_descent/tangents.gp}

\textbf{The disk at 73 for \(b_0=1\).} \S\ref{arith:63} gives
\(b\equiv1\pmod{73}\). At the three embeddings of \(K\) into
\(\mathbf Q_{73}\) the model is good and \(Y(P)\) is a unit.
Use the notation of \S\ref{arith:62}, with \(n_j=\#E_d(\mathbf F_{73})\) at
place \(j\). For the chart \(X=db\), the rows
\(L_j=\log([n_j]P_j)/73\bmod73\) and
\(V_j=n_jd/(2Y(P_j))\bmod73\) are
\[L=(30,34,22),\qquad V=(55,61,36).\]
Their pairwise determinants are \(33,16,28\) modulo 73. A global point
has logarithm proportional to that of \(P\) at all three places, with
the same rational proportionality factor. At least two embeddings
reduce to the same sign of \(P\). The corresponding determinant
function on that disk vanishes at 1 and its derivative is a unit.
Its derivative has integral power-series coefficients, so the
linear term dominates on the entire disk by the valuation estimate
of \S\ref{arith:62}. Thus \(b=1\). The B length-two chain centered at 1 shares
\(1/2\) with the fixed A chain. This excludes the class at every
rational height.
\cert{type1_descent/disk73.m}

\subsubsection{The discriminant-field cover}\label{arith:65}

Type 1 has no valid distinct-entry lift at any rational height.
In the two disks of \S\ref{arith:63} the only squared centers satisfying the
necessary center and singleton conditions are respectively 0 and 1.

\textbf{The discriminant character.} \S\ref{arith:64} handles the class of 1.
In the class of 0 write \(b=v^2\). If \(v\ne0\), then \(k=v_2(v)\ge3\).
In fact \(k<0\) would make \(5+4/v^2\) a square congruent to 5 modulo 8.
At the ramified prime of the cubic field of \S\ref{arith:64}, the normalized
valuation of its root \(r\) is \(-4\), and \(1-1/r\) is nonsquare.
For odd \(v\), substituting \(v^2\) for 1 does not change this local
squareclass, as \(v^2-1\in8\mathbf Z_2\). This contradicts \S\ref{arith:63}.
For \(k=1,2\), respectively, \((5v^2+4)/4\) is 6 or 5 modulo 8.

Use the sextic
\[h(B)=-800B^6+2625B^5-12050B^4+529B^3+54448B^2+53744B+14400.\]
The discriminant of the singleton cubic with leading coefficient 64 is
\(2^{26}h(b)\), so \(h(b)\) is a square. The sextic is irreducible.
Let \(F=\mathbf Q(\theta)\) with \(h(\theta)=0\). We claim that
\(s=1-b/\theta\) is square in \(F\).

Its norm is \(h(b)/14400\). Away from the primes above
\(D=\{2,3,5,23,41,2111,\allowbreak 11801,144461\}\), the polynomial has unit leading
and constant coefficients and unit discriminant. For integral \(b\), at
most one prime over a given rational prime divides \(b-\theta\), and it
has residue degree one. The square norm forces that valuation to be even.
At a pole the valuation is even because \(b=v^2\). This puts \(s\) in
the group of squareclasses unramified outside \(D\).

The discriminant of \(F\) is 14755078257381100. PARI certifies that its
factor base generates the class group, and all its generators are
principal. For each set of places used in this proof, the verifier
checks the expected number of independent unit squareclasses by
valuations and local squareclass maps. Thus all the global upper
spaces used below are unconditional.

Start with the norm-square subspace. The class of \(s\) is trivial at
every place above 2 and 73. At 2 one has
\(v_{\mathfrak p}(b/\theta)>2v_{\mathfrak p}(2)\) because \(k\ge3\);
at 73 the numerator \(b\) reduces to zero and \(\theta\) is a unit.
These conditions leave dimension eight. For each odd \(l\in D\), put
\(t_i=v_i(\theta)\) at the primes above \(l\). For integral \(b\), the
valuation vector of \(s\) modulo two is \((\max(t_i,0))_i\), up to
corrections supported where \(b\) and the integral \(\theta\) have the
same residue. Such a residue must be rational and must pass the two
center-square tests, the product-square test and the split square-root
test for the singleton cubic. For a pole at \(l\ne5\) the vector is
\((t_i)_i\) modulo two. A pole at 5 would give \(5+4/v^2\) of valuation one.
Taking the linear span of all these possibilities suffices. The
successive dimensions at \(3,5,23,41,2111,11801,144461\) are
\(6,5,4,2,1,0,0\). This proves the claim.
\cert{type1_descent/sextic.gp}

\textbf{The sextic elliptic cover.} Put \(d=-5\theta\) and
\[E:\quad Y^2=X(X-d\theta)(X+4d/5),\qquad X=db .\]
The three square conditions on \(b,5b+4,s\) give a point \(R\in E(F)\).
The split two-descent of \S\ref{arith:64} gives Selmer dimension five and hence
rank at most three. The five explicit basis triples are indexed by bits
of weights \(1,2,4,8,16\). The three nonidentity two-torsion points have
masks 27, 7 and 28. Their classes are distinct, so they span the image
of all torsion modulo doubling.
\cert{type1_descent/sextic_tangents.gp}

Use the completions \(F_1,F_2\) of ramification and residue degrees
\((2,1),(1,3)\) above 2. At each place scale the abscissa by \(c^2\)
and ordinate by \(c^3\), choosing the positive rational power of 2 for
which \(c\theta\) is a unit. Both models are integral. Write
\(m=105\cdot2^7\) and \(\Lambda_i=m\log E(F_i)\) on the scaled models.
The local abscissas in the data file give Kummer bases of dimensions
4 and 5. Hence their scaled logarithms generate the lattices.
Locally number these points starting at one. The logarithms of points
2,3 in the first list and 3,4,5 in the second form lattice bases.
The matrices of all the local basis logarithms modulo \(2\Lambda_i\) are
\[U_1=\begin{pmatrix}0&1&0&0\\0&0&1&0\end{pmatrix},
 \qquad U_2=(0_{3\times2}\mid I_3).\]
For the computation use \(z=-X/Y\) of the \(m\)-multiple. The invariant
differential is an integral even power series in \(z\), with constant
term one, so \(\log(z)-z\in z^3\mathcal O_{F_i}\). All nonzero
parameters computed are in \(2^7\mathcal O_{F_i}\). Multiplying either
inverse basis matrix by 256 makes it integral. Thus the indicated
bits, including those after one more division by 2, follow from the
precision in the certificate.
\cert{type1_descent/sextic_local_points.txt},
\cert{type1_descent/sextic_logs.m}

Let \(T=(0,0)\). The parameter of \(R-T\) is \(X/Y\) on the scaled
model and equals
\[\frac{\pm v}{2c\theta}
 \left((1-b/\theta)(1+5b/4)\right)^{-1/2}.\]
At both places the square-root correction is a unit congruent to
\(\pm1\) modulo 4, since \(k\ge3\) and \(v_2(\theta)\) is respectively
0 and 2. In the product of the two lattices, dividing
\(m\log(R-T)\) by \(2^{k-2}\) gives the primitive vector
\(u=(1,1,1,1,1)\) modulo two. Indeed the scaled leading term is an
odd rational unit times \(256/(c\theta)\), with either sign at each
place. The certificate checks its five bits. The cubic error in the
logarithm cannot change them since \(v_2(z)=k-1\) and the inverse
basis matrices have denominators dividing 256.

In particular \(R-T\) has infinite order. In the free group modulo
torsion write it as \(2^j Q\) with \(Q\) primitive at 2. The leading
vector of the product logarithm of \(Q\) must also be \(u\). On the
five-dimensional Selmer group the joint logarithm map modulo two has
rank two and misses \(u\). Its kernel consists of the four torsion
masks and their translates by 13. Thus the logarithm of \(Q\) is
even, and after a torsion translation its Kummer mask is 13.

\textbf{One more division.} Write \(e_i\) for the three roots of the elliptic
cubic and \(d_i\) for the entries of the mask-13 triple. Its projective
two-cover \(C\) has equations \(d_i z_i^2=X-e_i\) on the affine chart.
For each \(i\), omit \(z_i\) and take the tangent \(L_i\) to the resulting
conic at the supplied field point. Use the two squareclass functions
\(F_C=(L_2L_3,L_1L_3)\). Their values are unchanged modulo squares
by projective rescaling. They extend as squareclasses at a zero by
removing the even order of vanishing.

Here is the local comparison used in the calculation. If two points
on \(C\) differ by \(V\) in its Jacobian \(E\), the quotient of their
\(F_C\)-values is the ordinary Kummer class of \(V\). Indeed the
tangent section on \(C\) cuts out \(2D_i\), where \(D_i\) is a fiber
of the degree-two projection to the conic. The class \(D_j-D_k\) is
the two-torsion point \((e_i,0)\) for the complementary index \(i\):
the two deck involutions flip signs \(z_j,z_k\), and their product is
the deck translation of the two-cover fixing the square root of
\(X-e_i\). Thus \(L_j/L_k\), equal to \(L_jL_k\) modulo squares,
has divisor twice a representative of that two-torsion class.
Evaluation on a degree-zero divisor gives its usual Kummer coordinate.
On the target elliptic curve the images of the two points differ by
\(2V\).
\href{https://doi.org/10.1112/jlms/jds063}{Descent functions on coverings}

Take a local reference lift on \(C\) at each place in the calculation.
Write \(F_p\) for its pair of squareclasses and \(M_p\) for the complete
local Kummer image. Any global point on \(C\) gives a single global
pair \(g\) with \(g_p-F_p\in M_p\) in additive squareclass notation.
The support consists of the model's bad places, the primes of the
entries \(d_i\) and the primes of the coefficient ideals of the
tangents. Away from these places the primitive tangent has even
valuation on the smooth conic, as in \S\ref{arith:64}, and both functions have
even valuation. The verifier ranges over the entire resulting pair of
global squareclass spaces. It imposes the affine linear constraints
at each finite place in the support and at both real places.

At our two dyadic places the final abscissa in each data list fixes
the reference elliptic image \(P_i\). In the corresponding bases
\(m\log P_i/2\) has bits \((1,0)\) and \((1,1,1)\).
For every compatible global pair add to these bits the local images
under \(U_i\) of \(g_p-F_p\). The result is exactly the residue of
\(m\log Q/2\) for any global lift. In every solution of the linear
system the first coordinate at the ramified place plus the second
coordinate at the unramified cubic place is \textbf{one}. This functional
vanishes on zero and on \(u\). The logarithm of \(Q\) can therefore
neither be divisible by four nor have the required leading vector.
This contradiction rules out every nonzero \(v\) in the zero disk.
At \(v=0\) the reflected B chains coincide. Together with \S\ref{arith:64}
this proves the exclusion.
\cert{type1_descent/sextic_four.gp}

\subsection{Type 6. Descent on two Prym twists}\label{detail:type6}

\subsubsection{The three classes on the base cubic}\label{arith:66}

Give the cubic \(H_6(a,c)=0\) of \S\ref{app:type6} the origin
\(O=(1,0)\). It has rank two and torsion of order two. In its group
modulo 216 every valid type-6 configuration belongs to one of the
three classes represented by \((1,1),(9,1),(16,4)\).

\emph{Proof.} On a valid configuration the coefficient \(E_6\) in \S\ref{app:type6}
is nonzero. Reduction of the third moment equation modulo \(H_6\)
simplifies the B singleton square and the A singleton sum and product to
\[b=(5a-c+5)/4,\quad S=(-3a+15c+22)/4,\quad
 T=(8a^2-40ac-67a+80c^2+95c+149)/16.\]
Put \(c=v^2\). For \(v\ne0\) the substitution \(w=(a-4)/v\) yields
the smooth plane quartic with affine equation
\[vw^3-5v^2w^2+10v^3w-10v^4+3w^2-25vw+50v^2-40=0.\]
Its inverse is \(a=4+vw,\ c=v^2\). This gives a genus-three model for
the required square cover. In particular, at an odd prime \(p\nmid30\)
an integral candidate reducing to \((4,0)\) requires 40/3 to be a
square modulo \(p\). Indeed \(H_6(a,0)=(a-4)^2(a-1)\) and the
coefficient of \(c\) is \(-40\) at \((4,0)\). It follows that
\(w\) is a unit there and reduces to a root of \(3w^2-40\).

The minimal elliptic model is
\[y^2+xy+y=x^3+x^2-2913x+80031.\]
The rank bounds both equal two and the torsion subgroup has order two.
Put \(P=(4,0),Q=(4,1)\) on the plane cubic and let \(U=(1,1)-P\).
The point \(U\) has order two. Exact rational division verifies that
no nonzero binary combination of \(P,Q,U\) is divisible by two and
no nonzero combination of \(P,Q\) with coefficients in
\(\{0,1,2\}\) is divisible by three. Hence \(P,Q\) give bases modulo
9 and 27. Modulo 216 the classes are parametrized uniquely by
\(iP+jQ+tU\) with \(i,j\bmod216\) and \(t\bmod2\).

At each good prime below, enumerate the projective plane cubic.
Keep all points at infinity. For an affine point, test the three
center squares, the two singleton square roots of \(X^2-SX+T\)
with multiplicity, and the quartic condition at \((4,0)\).
Reduce the surviving points in the finite elliptic group modulo the
given multiple. Intersect with the reductions of the coefficient
classes, lifting the surviving classes between stages.

{
\begin{longtable}[]{@{}lll@{}}
\toprule\noalign{}
Multiple & Primes & Global classes left \\
\midrule\noalign{}
\endhead
\bottomrule\noalign{}
\endlastfoot
9 & 7, 11, 13, 17, 41, 61, 73, 307 & 4 of 81 \\
27 & 197, 271, 397, 647 & 3 of 729 \\
216 & 7, 13, 23, 31, 41, 43, 67, 73 & 3 of 93,312 \\
\end{longtable}
}

The final classes are \(P+U,P-Q+U,-P+2Q+U\). The group law gives the
representatives asserted in the statement. At those three points
\(c=1\) or 4, so the length-four B chain meets the centered A chain.
The next two arguments impose the covering conditions on all points in these cosets.
\cert{check_type6_classes.m}

\subsubsection{Two disks on a common Prym twist}\label{arith:67}

In the classes of \((1,1)\) and \((16,4)\) of
\S\ref{arith:66}, the only points of the base cubic with \(c\) a rational
square are the two class representatives, respectively. Both are
collisions. A valid type-6 configuration must lie in the class of
\((9,1)\).

\textbf{A common twist.} With origin \((4,0)\), the cubic has Weierstrass
coordinates \(t=(a-1)/c,\ Y=2A(t)c+B(t)\), where
\[A=t^3-5t^2+10t-10,\qquad B=-6t^2+5t+20.\]
Substitution gives \(Ac^2+Bc+9t-10=0\) and
\(Y^2=5t(32t^2-155t+192)\). Write \(K=\mathbf Q(r)\),
\(r^2=-551\), \(e=(155+r)/64\), and \(h=-B(e)/(2A(e))\).
For a representative \(R_0\), the two-torsion character
\((t(R)-e)/(t(R_0)-e)\) is a square in \(K\) throughout
\(R_0+216E(\mathbf Q)\). Changing the origin translates both points
and preserves their difference. With \(c=v^2\), the function
\(q_1=a-1-ec=c(t-e)\) therefore has constant squareclass.
For representatives 1, 2, 3 in the order of \S\ref{arith:66} we may use
\(\delta=-e,8-e,-e\), respectively. At \((16,4)\),
\((15-4e)/(-e)=((5+r)/16)^2\).

On the quartic of \S\ref{arith:66} in coordinates \((v,w,z)\), put
\[q_1=3z^2+vw-ev^2,\quad q_2=v^2-hz^2,\qquad
 q_3=\frac{q_2^2-H_6(4z^2+vw,v^2,z^2)/(A(e)v^2)}{q_1}.\]
Both divisions are exact. The quartic is \(q_1q_3=q_2^2\), and a
candidate lifts to \(R^2=\delta q_1, RS=\delta q_2,S^2=\delta q_3\).
If \(M_i\) represents \(q_i\), let
\(M(X)=M_1+2XM_2+X^2M_3\), \(g=-\det M(X)\). The sextic is
squarefree. Bruin's construction gives the genus-two Jacobian
\(J_\delta=\operatorname{Jac}(y^2=\delta g(X))\).
\cite[Theorem~5.1 and \S6]{BruinPrym}

For clarity, the divisors used below can be constructed without
choosing a point at infinity on the genus-two curve. Write
\(Z=(v,w,1), W=(1,w',0), L=R+XS\), differentiating with respect
to \(v\). Take the monic associate \(U\) of
\(\delta W M(X) W^t-(L')^2\), and reduce the second coordinate of
\((L'Z-LW)M(X)\) modulo \(U\) to obtain \(V\).
Then \(V^2=\delta g\pmod U\), by the tangent identity in the Type~13 construction of \S\ref{app:type13-details}
applied to \(\delta M(X)\) and division of the ordinate by \(\delta\).
The two support points minus a hyperelliptic fiber give a degree-zero
class, including when the support is nonsplit.

\textbf{Unconditional descent.} The quadratic and quartic factors of \(g\)
over \(K\) give absolute number fields of degrees four and eight.
They are isomorphic to the fields defined by
\[x^4-x^3+63x^2+313x+1544,\]\[
 x^8-3x^7+4x^6-328x^5+3937x^4-26339x^3+91290x^2-140562x+86400.\]
PARI certifies generation of the ideal class groups by the factor
bases with norms at most 5 and 83. Magma checks the field
isomorphisms, uses these bounds for the two fields and Minkowski
bounds for other fields, and obtains the rank upper bound 2 for
both twists by 2-descent.

\textbf{The disks at 197.} This prime splits in \(K\), and the reduced base
elliptic group has invariants \((2,108)\). Thus every point in each
global class specializes to its designated representative. Choose
\(v\equiv v_0\) with \((v_0,w_0)=(1,-3),(1,5),(2,6)\), respectively.
There is a single global sign of \(R\) for which the lift specializes
at \emph{both} places to the chosen
\(R_0^*=e,8-e,(7+5r)/32\), respectively. Indeed write the base point as
\(R_0+2D\) with \(D\in108E(\mathbf Q)\), so \(D\) reduces to zero.
The function
\[F(D)=\frac{t(R_0+2D)-e}{t(R_0)-e}\]
is a square of a rational function \(H(D)\) normalized by \(H(0)=1\).
The divisor of \(t-e\) is twice the divisor of the corresponding
two-torsion point minus the origin. Translation preserves this class
in \(\operatorname{Pic}^0\) and pullback by \([2]\) kills it. The
remaining constant is fixed by evaluation at zero. At each prime
above 197, \(F\) is a unit near the specialization of zero, and on
the smooth integral model the local ring is normal. Thus \(H\) and
its inverse are regular there and \(H(D)\equiv1\). Taking
\(R=R_0^*(v/v_0)H(D)\) gives the stated sign.

For the two places in the order \(r\equiv146,51\), use
\(\log([n_j]D)/197\) on each Jacobian block, where \(n_j\) is the
order of the reduced Jacobian. Use \(dX/(2y),X\,dX/(2y)\) as basis.
On the common twist the orders are 38196 and 42540. The collision
logarithms and the derivatives with respect to \((v-v_0)/197\)
reduce as follows.

{
\begin{longtable}[]{@{}lll@{}}
\toprule\noalign{}
Center & Collision logarithm & Derivative \\
\midrule\noalign{}
\endhead
\bottomrule\noalign{}
\endlastfoot
\((1,1)\) & \((123,77,4,92)\) & \((66,116,10,167)\) \\
\((16,4)\) & \((75,54,2,128)\) & \((148,119,66,137)\) \\
\end{longtable}
}

The verifier computes the logarithms by translating the multiple
to a split degree-two divisor and displacing its two abscissas modulo
\(197^2\). For the derivative write \(\dot U=dU/dv\).
Implicit differentiation at each root of \(U\) and interpolation give
\(-n_j[X](X^k\dot U/(2V)\bmod U)\) for \(k=0,1\).
All denominators in this chart are units throughout the disk. In
particular the two roots of \(U\) are distinct in the residue
algebra and their ordinates are units.

The determinants of the first three columns of the two logarithms
followed by the respective derivative are 7 and 100 modulo 197.
Hence the two global classes span the rational Mordell--Weil group
up to torsion, and on each disk an integral linear functional
annihilates their exact logarithms and is a unit on the tangent.
Its integral vanishes at the collision and at any rational lift.
The pulled-back differential has integral power-series coefficients
in \(v-v_0\). For \(m\ge2\) and \(k=v_{197}(v-v_0)\ge1\),
\(mk-v_{197}(m)>k\), so the linear term of the integral dominates
every higher term. Consequently \(v=v_0\); the quartic has a unique
\(w\) in its disk since its \(w\)-derivative at these centers is a
unit. This proves the proposition.

For the twist \(\delta=8-e\) treated in \S\ref{arith:68}, the same calculation gives
orders 39780 and 42540, collision logarithm \((184,142,27,148)\),
and tangent \((69,56,174,18)\). The logarithm proves that the
collision divisor has infinite order. The certified rank interval
for this twist is \([1,2]\).

\cert{type6_descent/class_bounds.gp},
\cert{type6_descent/prym.m}

\subsubsection{The remaining dyadic disk}\label{arith:68}

Type 6 has no valid distinct-entry chain at any rational
height. In the class of \((9,1)\) the full necessary chain conditions force
\((a,c)=(9,1)\).

\textbf{The base disk.} Use the generators and origin of \S\ref{arith:66}.
None of the 48 nonzero combinations of the generators modulo 7 is
divisible by 7. Of the 49 lifts of the class of \((9,1)\) from 216 to
1512, the reduction tests of \S\ref{arith:66} leave seven at 67 and one after 79.
It is the class of \((9,1)\) itself. On the minimal elliptic model,
the multiples by 1512 of both generators lie in the fifth formal
subgroup at 2, as checked with the parameter \(-x/y\).
The finite quotient by this subgroup has order divisible only by 2
and 7: the reduction is split multiplicative of type \(I_{14}\).
The saturation checks at 2 and 7 therefore give the same inclusion
for the full rational Mordell--Weil group.
\cert{check_type6_classes.m}

In the Weierstrass coordinates of \S\ref{arith:67} the minimal coordinates satisfy
\(x=40t-65,\ 2y+x+1=40Y\). Hence the invariant differential is
\[\frac{dx}{2y+x+1}=\frac{dt}{Y}=-\frac{dc}{(H_6)_a}.\]
The last identity follows by differentiating \(A c^2+B c+9t-10=0\).
At \((9,1)\), the partial derivatives of \(H_6\) are 40 and \(-120\).
Putting \(c=1+64z,\ a=9+64u\) gives, after division by \(64\cdot40\),
an equation \(u-3z+8G(z,u)=0\) with integral higher terms.
It has a unique integral analytic branch through the origin and
\(40/(H_6)_a\) has integral series with constant 1 there. For
\(c=1+256s\), the elliptic logarithm divided by \(-32/5\) is therefore
\(s\) plus a strictly contracting integral series: each coefficient
of degree \(j+1\ge2\) has valuation at least \(2j-v_2(j+1)\).
This parametrizes the whole translate of the fifth formal subgroup,
by the elliptic formal exponential and the inverse function theorem.
Indeed the parametrization and the formal exponential of the integral
agree near zero and then on the disk by analytic continuation.
Consequently \(c\equiv1, a\equiv9\pmod{256}\). Choose the rational
sign of \(v=\sqrt c\) with \(v\equiv1\pmod{128}\); then
\(w=(a-4)/v\equiv5\pmod{128}\).

\textbf{The possible second global direction.} Take \(d=8-e\) and use the
reciprocal model \(y^2=f(X)=X^6 d g(1/X)\) of \S\ref{arith:67}. Write \(J\)
for its Jacobian, \(D\) for its collision divisor, \(T\) for its
nonzero rational two-torsion point, and \(F\) for a hyperelliptic
fiber. Write \(\theta\) for the class of \(X\) in the sextic algebra.
For a divisor of odd degree one can evaluate \(X-\theta\) modulo
squares and base-field scalars. The norm of this class equals the
leading coefficient of \(f\) modulo squares. Differences of two such
classes belong to the usual even descent space.
\href{https://math.mit.edu/~poonen/papers/descent.pdf}{Poonen and Schaefer, §§5, 11--13}

Here is the affine descent calculation. Put \(S\) equal to the ten
places over \(2,3,5,7,31\). The monic sextic is integral, its
discriminant a unit, and the leading coefficient of \(f\) a unit
outside \(S\). At such a place an integral abscissa specializing
to a root in the étale algebra has even order in that component
because the branch point is simple and \(y^2=f(X)\). At a pole all
components have the same valuation. Applying this to closed points
over their residue fields and taking norms gives diagonal valuation
parity for arbitrary divisors. The class group of \(K\) is cyclic
of order 26 and a dyadic prime generates it. Thus a base-field
scalar clears these diagonal parities outside \(S\). The two factor
fields of \S\ref{arith:67} give a space of dimension 40 of squareclasses with
even valuations outside \(S\), before taking the quotient by scalars.

At each listed place the verifier takes an odd local point as origin.
It adjoins the scalar classes and the differences of odd points, the
class of \(D\), and the class of \(T\). At one dyadic place an additional
degree-two divisor has support \(X^2+5/2\); extracting a square in
\(\mathbf Q_2(\sqrt{-10})\) verifies its ordinate. The local spaces
have dimensions 8 and 6 at 2, 4 and 4 at 3, 3 and 3 at 5, 2 and 4
at 7, and 4 and 2 at 31, in the order of the certificate.

These spaces are complete. With \(n\) irreducible factors at an odd
place the dimension including scalars is \(n\), and at the two dyadic
places it is \(n+3\). In fact, when there is an odd factor the scalar
map is injective, rational two-torsion has dimension \(n-2\), and
the local descent is injective modulo doubling. At the two remaining
odd places the degrees are \((2,4)\). There the two-torsion dimension
is one. If both fields contain the same quadratic extension then the
scalar map loses one dimension and the roots admit a conjugate
three-plus-three partition. Otherwise the scalar map is injective
and there is no such partition. The kernel criterion of
Poonen--Schaefer gives total dimension two in either case.
All ten places have local points; at all other finite places good
reduction gives an odd-degree divisor by passing to a sufficiently
large odd unramified extension. Thus any global rational divisor
class of degree one is represented by a divisor, by the local
Brauer obstructions and the Brauer--Hasse--Noether theorem.

The norm and local affine constraints, homogenized with one parity
coordinate, have a 14-dimensional nullspace. There is an odd vector,
so the even subspace has dimension 13. Scalar classes account for 11:
the quartic relative factor field has no quadratic intermediate field
(the absolute octic field has just one quadratic and no quartic
subfields). The images of \(D,T\) are independent, as the local
coordinates below show. They therefore span the even space modulo
scalars.

At the two dyadic places in the order \(r\equiv155,101\pmod{256}\),
choose a local Weierstrass point \(W_0\) as in the certificate. The
free local coordinates use the differences from \(W_0\) of the points
with abscissas 352 and 4704 at the first place. At the second use
the point at 9 minus \(W_0\), and the quadratic divisor at
\(X^2+5/2\) minus a fiber. Together with three and one independent
Weierstrass differences, respectively, these give bases modulo
doubling. In particular the two-primary local torsion has exponent
two. Sign choices of these free generators give the same coordinates
modulo two. For local torsion coordinates followed by free coordinates,
\(D\) has rows \((0,1,1;0,0)\), \((0;0,0)\).
The rows of \(T\) are \((0,1,1;0,0)\), \((1;0,0)\).
Every possible odd class \(W\) has
\[\log(W-W_0)\equiv(1,1,0,1)\pmod {2\Lambda},\]
where \(\Lambda\) is the product of the free local logarithm lattices.
Indeed the verifier takes any odd solution of the affine system,
subtracts the local Weierstrass images and reads the two free
coordinates at each place. Adding either even generator preserves them.
\cert{type6_descent/beta_selmer.m}

If the rank is two, \(J(K)/2J(K)\) has dimension three. Its even
descent image has dimension two, and Poonen--Schaefer's Theorem 11.3
writes the extra kernel generator as \(B=2W-F\) for a degree-one
class, modulo doubling. Choose that representative. Since \(F=2W_0\)
locally, \(\log(B)/2\) gives the displayed row. The classes
\(D,T,B\) span modulo doubling and \(D,B\) span up to finite index
and torsion. In rank one, \(D\) alone spans up to finite index and
torsion.

\textbf{The disk calculation.} In the same free local coordinates the
collision logarithm divided by two and the derivative at \(v=1\)
multiplied by eight have respective reductions
\[(1,1,1,0),\qquad (1,0,0,0).\]
The first two proposed global rows \((1,1,1,0),(1,1,0,1)\) are
independent modulo two. Thus even after saturation the reduction of
the global logarithm space is contained in their span, also in the
rank-one case.

For details of precision, the verifier multiplies each divisor by
\(2^{18}\cdot315\) and translates it to two split points with ordinates
\(y_j\). If their abscissa displacements are \(h_j\), integrate in
the variable \(X=x_j+h_j z\). The degree-\(k\) coefficients of
\(f(x_j+h_j z)/y_j^2-1\) have valuations at least \(6k\) for
\(1\le k\le6\). The coefficients of the inverse square root of \(f(x_j+h_j z)/y_j^2\)
with constant 1 have valuations at least \(4k\). After division by
the integration index and the multiplier, the discarded terms of
degrees eight and higher give error in \(2^8\mathbf Z_2\) in both
logarithms. The inverse basis matrices have entry valuations at
least \(-1\), so this precision determines the displayed reductions.

On the reciprocal curve the two pulled-back differentials are
\[\gamma(R,-S)\frac{dv}{F_w},\qquad
 \gamma=(58773-2161r)/128000,\]
where \(F(v,w)=0\) is the quartic and \(R^2=d q_1, RS=d q_2\).
Indeed the canonical complete intersection of three quadrics has
canonical bundle \(\mathcal O(1)\). Its two anti-invariant
differentials have basis \(R\,dv/F_w,S\,dv/F_w\), obtained by
multiplying the plane-quartic differential by the extra linear
coordinates. Changing the sheet negates the Mumford image, so the
two pullbacks lie in this space. Their constant and linear jets at
the collision determine the coefficients: the basis jet determinant
is nonzero and the verifier evaluates both using the monic tangent
support.

Here is a uniform bound for integration throughout the required disk.
Put \(v=1+32z, w=5+32u\). The scaled quartic has unit \(u\)-derivative
and gives an integral series \(u=z+4z^2\mathbf Z_2[[z]]\).
Also \(40/F_w\) is integral with constant one. Write \(L_j\) for
the local basis matrix. The vectors \(\gamma d(1,0)L_j^{-1}\),
\(\gamma(1-h)(0,1)L_j^{-1}\), \(\gamma(0,1)L_j^{-1}\) have entry
valuations respectively \((-5,-2),(-5,-2),(-5,-2)\) at the first
place and \((1,2),(4,3),(0,-1)\) at the second. On the positive
sheet \(R/d=1\) at the origin. At the first place
\(v_2(e)=v_2(d)=-5\) and \(q_1/d-1\in64z\mathbf Z_2[[z]]\).
Thus \(R/d-1\) and \(S-(1-h)\) belong to \(32z\mathbf Z_2[[z]]\).
At the second place \(v_2(e)=6,v_2(d)=3,v_2(1-h)=4\);
\(q_1/d-1\in8z\mathbf Z_2[[z]]\) and \(S=(v^2-h)/(R/d)\).
Together with integrality of the constant row, these bounds make
eight times the differential coefficient integral in the basis on
both disks. These parametrizations on the open disk \(|z|_2<1\)
have analytic coordinates and the parameter is recovered from \(v\).
On such a disk the abelian logarithm displacement equals the
termwise primitive, by functoriality from the Abel--Jacobi map.
\href{https://www.mathe2.uni-bayreuth.de/stoll/papers/uniform-bounds-2015-11-25.pdf}{Stoll, §7, integration on an open disk}

For \(k=v_2(v-1)\ge7\) we have \(v_2(z)\ge2\). After
integration and division by \((v-1)/8\), each nonconstant term vanishes
modulo two since \(2j-v_2(j+1)\ge1\) for \(j\ge1\). The displacement
therefore has primitive row \((1,0,0,0)\), outside the global span.

Finally the two local signs of \(R\) need not be prescribed. Negate
the global divisor if necessary to obtain the positive sign at the
first place. If the sign is negative at the second, subtraction of
the global collision logarithm adds the second block of \(-2\log D\).
After division by four this has row \((0,0,1,0)\), also outside the
global span, while the varying displacement is divisible by 16.
If the signs agree the preceding tangent test rules out every
\(v\ne1\). At \(v=1\) the unique local quartic branch has \(w=5\),
so \(c=1,a=9\) and the chains collide. \S\ref{arith:67} handled the
other two classes.
\cert{type6_descent/beta_logs.m}

\subsection{Type 7. Descent from a rational branch point}\label{detail:type7}

\subsubsection{Five classes on the cubic}\label{arith:69}

With origin \((a,c)=(1,1/4)\), the type-7 cubic
of \S\ref{app:type7} has rank two and trivial torsion. A valid chain must be in
one of five classes modulo 216, represented by
\[(1,1/4),\quad(9,9/4),\quad(4,49/4),\quad(4,1/4),\quad(1,25/4).\]

\emph{Proof.} On \(H_7=0\) the linear subresultant of the two remaining
moment equations gives
\[b=(10a-4c+15)/6,\qquad S=(16a+20c+81)/12,\]\[
 T=-(128a^2-640ac+48a+80c^2-1380c-1701)/144 .\]
The resultant in \(a\) of its leading coefficient and \(H_7\) is
proportional to
\[(4c-9)^4(4c-1)^6(4c+11)
 (144c^2+152c+529)(144c^2+440c+1225).\]
The first two factors give collisions, the linear factor at negative
\(c\) is impossible for a square, and the two quadratics have negative
discriminants. The formulas therefore apply to every valid lift.
\cert{check_all_shift_chain_types.py}

The minimal elliptic model is \(y^2=x^3+x^2-13533x-323937\).
The two rank bounds agree. With
\(P=(-7/2,1/4),Q=(9,9/4)\) on the plane cubic, exact division at
2 and 3 proves that these points generate the quotient modulo 216.
At good primes apply the projective reduction test of \S\ref{arith:66} using
the three center squares and the two square roots of \(X^2-SX+T\).
The successive moduli 9, 27 and 216 leave seven, five and five
classes. At 216 they are \(0,Q,-P+Q,2P-2Q,3P-2Q\), respectively.
All five representatives are collisions. For \((4,49/4)\) one has
\(b=1\), so the two length-two chains share \(3/2\). For
\((1,25/4)\) one has \(b=0\), making the reflected B chains
coincide. The other representatives have \(c=1/4\) or \(9/4\).
Sections~\ref{arith:70} and~\ref{arith:71} determine the points satisfying the necessary chain conditions in all five classes.

The square cover \(b=v^2\) has a rational branch point at
\((1,25/4)\). Restriction to the zero line
\(a=(4c-15)/10\) gives
\[H_7((4c-15)/10,c)=-27(4c-25)(144c^2+440c+1225)/125.\]
The other two finite branch points are defined over
\(\mathbf Q(\sqrt{-5})\), and the three simple poles lie at infinity.
Riemann--Hurwitz gives genus four. The descent in \S\ref{arith:70} starts from the rational branch point.
\cert{check_type7_classes.m}

\subsubsection{Descent at a rational branch point}\label{arith:70}

In four of the five classes in \S\ref{arith:69}, the
square and splitting conditions for a type-7 chain force the collision
representative. The remaining class is that of \((9,9/4)\).

Write \(u=a-1\). The line \(a=1\) cuts the cubic in \((1,25/4)\) and
twice the origin \((1,1/4)\). On the genus-four cover of \S\ref{arith:69}, put
\(t=b/u\). Then \(c=(25+(10-6t)u)/4\), and the cubic divided by
\(72u\) reads
\[A u^2+B u+C=0,\qquad
 A=45t^3-145t^2+135t-27,\quad
 B=-8(45t^2-80t+21),\quad C=144(5t-2).\]
Putting \(y=(2Au+B)/8\) gives
\[y^2=5F(t),\qquad F(t)=27t^3-79t^2+57t-9.\]
The inverse recovers \(u\) generically, so this is an isomorphism
of the smooth projective base curves. Its only pole of \(t\) is the
origin. At the other point on \(a=1\), \(t=2/5\). The Weierstrass
coordinates \((135t,135y)\) have coefficients
\([0,-395,0,38475,-820125]\).

For any other lift, \(t=d z^2\), where \(d\) is a positive squarefree
divisor of 30. Indeed \(5F(t)<0\) for \(t\le0\). At a prime \(p>5\)
a pole of \(t\) has even order by the cubic equation. At a zero of
its reduction the leading and constant coefficients of the quadratic
in \(u\) are units, so \(u\) is a unit and \(tu=b\) is square.
To select \(d\), reduce the coset \(R_0+216E(\mathbf Q)\) at
\(11,17,19,41\). At a pole the unit squareclass is 15 by the leading
coefficient 135. At a zero use \(27u^2+168u+288=0\) and the
squareclass of \(u\). The five cosets give the following results in the
order of \S\ref{arith:69}. For the additional disk tests in the last
column impose all five square conditions of that proposition on every
finite fiber. A vanishing leading coefficient is retained.

{
\begin{longtable}[]{@{}lrrr@{}}
\toprule\noalign{}
\((a,c)\) & \(t_0\) & \(d\) & Prime forcing \(t\equiv t_0\) \\
\midrule\noalign{}
\endhead
\bottomrule\noalign{}
\endlastfoot
\((1,1/4)\) & \(\infty\) & 15 & \\
\((9,9/4)\) & 2 & 2 & 17 \\
\((4,49/4)\) & \(1/3\) & 3 & 17, 83 \\
\((4,1/4)\) & 3 & 3 & 61 \\
\((1,25/4)\) & \(2/5\) & 10 & 389 \\
\end{longtable}
}

\textbf{Cubic-field elliptic covers.} Let \(e\) be a root of \(F\) and write
\(F(t)=27(t-e)(t^2+g_1t+g_0)\). In \(K=\mathbf Q(e)\) the
two-descent character on the monic base model is \(135(t-e)\).
Every point in a fixed coset has the character of its representative
because 216 is even. Thus \(t-e\) has class \(\delta=t_0-e\) for a
finite representative and class 135 at the origin. Put \(k=135d\delta\).
Multiplying the square conditions gives the elliptic lift
\[X=kt,\qquad Y^2=X(X^2+kg_1X+k^2g_0).\]
In the five rows its ranks over \(K\) are respectively \(0,3,3,2,2\).
The first curve has torsion of order two, with abscissas zero and
infinity. Since \(5F(0)=-45\), only the origin remains.

The primes 61 and 389 are inert in \(K\). For the corresponding
rank-two curves take the two global points in the logarithm certificate.
At the anchor the ordinate is a unit. If \(N\) is the local reduction
group order, the coordinates of \(\log([N]P_i)/p\) modulo \(p\)
are obtained from the formal parameter \(-X/Y\) modulo \(p^2\).
The third row is the tangent \(k/(2Y_0)\). In the power basis of
\(h^3-8h-2=0\), the determinants, up to sign, are 31 at 61 and
216 at 389. Both reduction orders are prime to \(p\).
The first two logarithms thus span the global logarithm space over
\(\mathbf Q_p\), and the tangent lies outside it. There is a single
place, so negating the lift chooses the branch through the global
anchor. A linear form annihilating the first two rows and taking a
unit value on the third gives a function with a unique zero on the
whole disk. Indeed the invariant differential has integral coefficients
in \(t-t_0\), since its denominator is a unit. After integration a term
of degree \(j\ge2\) has valuation at least
\(jv_p(t-t_0)-v_p(j)>v_p(t-t_0)\).

\textbf{A second descent for the rank-three rows.} On
\(C_d:y^2=5F(dz^2)\), the two-cover calculation gives four possible
classes for \(d=2\) and eight for \(d=3\). For the rows with
\(t_0=2,1/3\), write \(z_0=1,1/3\), respectively. At 17, the
only admissible classes in the sextic algebra are represented by
\(z_0-\theta\) and \(-z_0-\theta\), where \(F(d\theta^2)=0\).
For a component of even degree test that the norm of the ratio is
square. For the odd-degree components the square characters must agree,
since the class is defined modulo a common rational scalar. The
certificate tests every class returned by two-cover descent. Changing
the global sign of \(z\) exchanges the two classes, so take
\(z_0-\theta\). This is the two-cover map of Bruin and Stoll.
\href{https://arxiv.org/abs/0803.2052}{Two-cover descent on hyperelliptic curves}

Put \(n^2=d/3\), \(s=nz\), and let \(r\) satisfy
\(9r^3-5r^2-3r+1=0\). Use \(n=1\) for \(d=3\). The factorization
\[F(3s^2)=9(9s^3-5s^2-3s+1)P(s),\qquad P(s)=9s^3+5s^2-3s-1\]
yields an elliptic quartic over \(L=\mathbf Q(r,n)\). In the chosen
two-cover class the product of the four factors corresponding to
\((s-r)P(s)\) has the squareclass of its value at \(s_0=nz_0\).
Indeed the four factors form a Galois-stable subset over \(L\) and
the common rational scalar occurs to the fourth power. With
\(k=(s_0-r)P(s_0)\), \(m=kP(r)\) the lift and its Weierstrass
coordinates are
\[q^2=k(s-r)P(s),\qquad X=m/(s-r),\quad Y=mq/(s-r)^2,\]
\[Y^2=X^3+kP'(r)X^2+km(27r+5)X+9km^2 .\]
For \(d=3\) its rank is one. At 83 the selected class forces
\(z\equiv1/3\): the component characters for the other sign disagree.
All three places have degree one. Two lift signs agree, and the minors
of the anchor logarithm and its slope for the three pairs of places
are nonzero (in cyclic order, 67, 70, 22). The same integral estimate
as above bounds each disk by the single anchor. This proves \(t=1/3\).

Finally the pairs of quadratic roots \(u\) at \(t=1/3,3,2/5\)
are \(\{3,-9/2\},\{3,13/6\},\{0,-760/167\}\). In each the second
value makes \(1+u\) a nonsquare. This proves the proposition.
\cert{type7_descent/classes.m},
\cert{type7_descent/covers.m},
\cert{type7_descent/odd_logs.gp}

\subsubsection{The final sextic cover}\label{arith:71}

Type 7 admits no valid distinct-entry chain at any
rational height. In its remaining class the necessary chain conditions
force \((a,c)=(9,9/4)\).

We use the local divisibility criterion of Selmer-group Chabauty.
Here the two-isogeny and a lift of one abscissa class determine the
local information it requires.
\href{https://arxiv.org/abs/1506.04286}{Stoll, \emph{Chabauty without the Mordell--Weil group}}

\textbf{Global descent on the sextic cover.} Use the last quartic of \S\ref{arith:70}
with \(n^2=2/3\) and \(s_0=n\). Shift its abscissa by
\(j=-m/(2r)\), writing \(x=X-j\). The point over \(s=-r\) is now
\(T=(0,0)\), and the curve and its two-isogenous curve have models
\[E:y^2=x(x^2+A x+B),\qquad E':y'^2=x'(x'^2-2A x'+D),
 \qquad D=A^2-4B .\]
Here \(A=kP'(r)+3j\), \(B=3j^2+2jkP'(r)+km(27r+5)\).
Write \(s_i\) for the four elements listed in the certificate, and
\(s_* = s_2\). The global image under the abscissa character
\(\alpha\) for \(E\) lies in \(\langle s_1,s_2,s_3,s_4\rangle\).
The isogeny Selmer space for \(E'\) is exactly \(\langle D\rangle\).
In particular a point of \(E(L)\) of trivial abscissa class belongs to
\(2E(L)\): it lifts along the dual isogeny to \(E'(L)\), where the
two possible classes are represented by the identity and the kernel
point of the dual isogeny.

Here are details of the finite computations. The field has model
\(h^6-96h^4+144h^2-54=0\), class number one and one place above 2,
of ramification index six. PARI certifies its class group and units.
On each curve the abscissa class is unramified outside the primes
dividing \(2BD\) (including denominator ideals). Form the squareclasses
of units with these primes inverted and intersect with the local
images at those primes and all real places. For the two local images
at a finite prime \(v\), the dimensions sum to
\(2+[L_v:\mathbf Q_2]\) above 2 and to 2 elsewhere. This follows by
composing the two isogenies and using the order of the local group
modulo doubling. The certificate finds local points until the paired
images attain this dimension. At a real place the negative abscissas
occur precisely when the model has a negative root and a nonempty
interval of negative abscissas with positive cubic value. The two
resulting dimensions are 4 and 1, giving rank at most three for the sextic cover. The certificate identifies the basis
by testing squares in the number field. Locally at 2 the kernel on
the four-dimensional space is exactly \(\langle s_*\rangle\).

\textbf{Separating the two classes in this kernel.} Let
\(e=(-A+\sqrt D)/2\). A global point of nonzero abscissa in the
class \(s_*\) can be written \(x=s_*u^2,\ y=s_*uv\), with
\(v^2=s_*u^4+A u^2+B/s_*\). Choose the global point \((U,V,1)\)
on \(s_*U^2+AUV+(B/s_*)V^2=1\) specified in the certificate.
Put \(\gamma=s_*U-eV\). The norm of \((x-e)/\gamma\) is \(v^2\).
Hilbert 90, or the identity \((w+v)^2/w=\operatorname{Tr}(w)+2v\)
when \(w\bar w=v^2\), shows that
\[x-e=\lambda\gamma\cdot\text{square in }L(\sqrt D),
 \qquad \lambda\in L^\times .\]
Where the tangent value is nonzero, its squareclass in the quadratic
extension is that of \(\lambda\):
\[\mathcal L(u,v)=(2s_*U+AV)u^2+AU+2BV/s_*-2v.\]
In a fixed local abscissa class the possible \(\lambda\)'s form one
coset of the local abscissa image of \(E'\). Indeed the second
two-descent character is \(x-e\). On the image of \(E'\) under the
dual isogeny it has class \(x'\), by the identity
\(x(\widehat\phi(P'))-e=(x'-\sqrt D)^2/(4x')\).
The kernel of squareclass extension is generated by \(D\), itself an
abscissa class of \(E'\). These identities also describe the characters
at the exceptional points by continuation of the usual descent map.

The certificate exhibits one element \(\kappa\) in every required local
coset. It evaluates \(\mathcal L\) on a local quartic point and tests
membership of \(\kappa\mathcal L\) in the local image of \(E'\).
At untested odd places the curves have good reduction, the coefficients
of the model are integral, and \(s_*,D,\kappa\) and \(N(\gamma)=s_*\) are units.
The conic point is integral there: the minimum valuation of its tangent
coefficients is zero and the conic matrix is invertible. Thus the
second character and \(\gamma\) are unramified and the local image on
\(E'\) consists of both unramified classes. The finite test includes
the support of the fractional ideal generated by the tangent coefficients.
The real signs are tested by exact number-field signs. Consequently the
global ratio \(\lambda/\kappa\) is in the dual Selmer space
\(\langle D\rangle\). Every global point in this class satisfies
\(x-e=\kappa\gamma\cdot\mathrm{square}\) in the quadratic extension.

Write \(M=L_v\) at the unique dyadic place, with valuation \(v(2)=6\).
Here \(s_*\) is square but \(\kappa\gamma\) is nonsquare in
\(M(\sqrt D)\). Both \(\kappa(2s_*U+AV\pm2\sqrt{s_*})\) are
nonsquares in \(M\), the trace-and-norm criterion for this assertion.
Hence \(E(L)\cap2E(M)=2E(L)\). The local two-torsion consists of
\(0,T\) because \(D\) is nonsquare. Since \(B=\alpha(T)\) is
nonsquare, \(T\) has no local half. On quotienting by \(\langle
T\rangle\), the local group has no two-torsion and the global subgroup
still has the same divisibility property. It applies repeatedly.

\textbf{The whole dyadic disk.} On the minimal rational base model the
component group at 2 has order 3 and its nonsingular reduction
has order 2 (additive reduction). The third formal subgroup has index
24, so contains \(216E_{\rm base}(\mathbf Q_2)\).
At the representative \((138,\pm675)\) the first abscissa derivative
in the formal group parameter is \(\pm1350\), and higher coefficients
are integral. As the minimal abscissa is \(135t-132\), this yields
\(t\equiv2\pmod {16}\), hence \(z\equiv\pm1\pmod4\).
For \(z\equiv-1\) the unit \((nz-r)P(nz)/k\) has the nonsquare
class of its value at \(-1\): changing \(z\) by a multiple of four
changes this unit by an element of valuation at least 15. Thus
\(z=1+4\xi\) with \(\xi\in\mathbf Z_2\).

Choose the sign of the lift so that \(q\) is on the branch through
\(q(n)=k\). Write \(H(\xi)=P(1+4\xi)-P(1)\) for the difference of the
elliptic lifts, using \(P(z)\) here for the point. For \(\xi\ne0\)
its formal parameter has valuation \(9+v(\xi)\). Indeed the
abscissa difference before subtraction has valuation \(25+v(\xi)\)
and the ordinate sum has valuation 16 by the quartic formulas. The
chord slope therefore has valuation \(-9-v(\xi)\).
On the seventh formal subgroup the logarithm and exponential give
inverse isometries with \(\mathfrak p^7\) for the displayed integral
model, since \(7>v(2)\). In particular \(\mathfrak p^{19}\) maps into
four times this subgroup.

Set \(g(X)=(n+4nX-r)P(n+4nX)/k\). Its constant is 1 and the higher
coefficients have valuations \(17,32,47,60\). Pulling back
\(dx/(2y)\) gives
\[\log H(\xi)=-\frac{2n}{k}\int_0^\xi g(X)^{-1/2}\,dX .\]
Its linear and quadratic coefficients have valuations 9 and 14.
All coefficients of higher degree have valuation at least 19. For
degrees 3 through 12 the valuations are
\(25,24,35,40,51,44,55,60,71,70\). For degree \(i+1\), \(i\ge12\),
the binomial expansion gives the lower bound \(9+3i-6v_2(i+1)\):
a degree-\(i\) product contributes at least \(15i\), and its binomial
denominator costs at most \(12i\). This bound is at least 19 and tends
to infinity.

Write \(\xi=2^j u\) with \(u\) an odd dyadic unit. Modulo
\(\mathfrak p^{19}\), \(2^{-j}\log H(\xi)\) equals
\(u\log H(1)\) for \(j=0\) and \(u\log H(2)/2\) for \(j\ge1\).
For odd \(u\) the quadratic difference has the extra factor \(u^2-u\);
for \(j\ge1\) the normalized quadratic term already has valuation at
least 20. Modulo four times the local group, \(u\) can be replaced by
an odd integer in its residue class modulo four.

It remains to check two points. Modulo \(T\), \(H(1)\) admits one
successive local halving and \(H(2)\) admits two. In both cases choose
the final representative \(R\) modulo \(T\) having square abscissa.
Neither \(R\) nor \(R+T\) has a local half. Moreover both
\(\kappa\mathcal L(u,v)\) and \(D\kappa\mathcal L(u,v)\) for this
representative are nonsquare, using \(u^2=x(R)/s_*\). Thus its class
modulo doubling cannot come from a global point. Only the global
abscissa classes 1 and \(s_*\) could give its local abscissa;
the first would be globally divisible, and the second fails the
tangent test. Translation by \(T\) preserves the image of the global
subgroup. The certificate checks halving by extracting square roots
of \(x\) and \(2x+A-2y/\sqrt x\), trying both signs, and checks the
precision of each square test. By the logarithm congruence every
\(H(\xi)\) with \(\xi\ne0\) has this same obstruction after division
modulo \(T\): division is unique there and multiplication by an odd
integer acts trivially modulo two. Global divisibility proved above
would require the primitive class to be global. Hence \(\xi=0\).

This forces \(t=2\). The quadratic gives \(u=a-1=8\) or \(144/23\).
The latter makes \(a=167/23\) a nonsquare. The former gives
\((a,c)=(9,9/4)\), a collision by \S\ref{app:type7}. \S\ref{arith:70} treats the
other four classes.
\cert{type7_descent/sextic.gp},
\cert{type7_descent/dyadic.m}

\subsection{Type 8. Reciprocal elliptic covers}\label{detail:type8}

\subsubsection{Product curves and three possible squareclasses}\label{arith:72}

Let \(b=v^2\) be the squared B center in a valid
type-8 chain. Every odd prime in the reduced denominator of \(v\) is
1 modulo 20, and \(v_2(v)\ge-1\). Put
\[G(b)=5(256b^4-1152b^3+1824b^2-388b+405).\]
The curve \(Y^2=G(b)\) has rational rank five, and the Jacobian of
\(C:Y^2=bG(b)\) has rational rank four. In the two-cover descent on
\(C\), each valid chain belongs to the class represented by
\(b_0-\theta\) for one of \(b_0=1/4,9/4,25/4\), where
\(\theta G(\theta)=0\) in the associated étale algebra.

\emph{Proof.} For the B polynomial put
\(B_v(X)=\prod_{j=-2}^2(X-(v+j)^2)\). The singleton polynomial of
\S\ref{app:type8} satisfies
\[(X-1/4)K_b(X)=B_v(X)-B_v(1/4),\qquad 256K_b(0)=G(b).\]
All four roots of \(K_b\) are rational squares. Their product gives
the quartic model; multiplication by the square \(b\) gives the
genus-two model. A vanishing B center already repeats vertices.

Suppose \(p\) is an odd denominator prime, or \(p=2\) and
\(v_2(v)<-1\). After substituting \(X=bU\), the monic quartic
\(K_b(bU)/b^4\) is integral locally. Its four roots are integral
squares. The five normalized B vertices have square tending to 1
modulo \(p\), and the normalized central A square is zero modulo
\(p\). The reduction is therefore \(((U-1)^5+1)/U\). At 2 this
polynomial has no linear root. At 5 it is \(U^4\); all four roots
would be divisible by 5, while the constant term \(G(b)/(256b^4)\)
has valuation exactly one. At every other odd prime the four roots
are \(1-\zeta^i\), for \(1\le i\le4\) and a primitive fifth root of
unity \(\zeta\). Splitting requires \(5\mid p-1\). The ratio
\((1-\zeta^{-1})/(1-\zeta)=-\zeta^{-1}\) is square, and \(\zeta\)
is square because its order is odd. Thus \(-1\) is square as well,
giving \(p\equiv1\pmod {20}\).

The quartic with the point over \(b=1/4\) has minimal model
\(y^2=x^3-105150x+7613125\), with both rank bounds equal to five.
For the genus-two Jacobian the rank upper bound is four and the
rational torsion has order two. The points over
\(1/4,3/4,9/4,4,25/4\), based at the rational point at infinity,
give five independent images under the two-descent map. This proves
the matching rank lower bound.

Two-cover descent gives seven possible classes on \(C\). At
\(13,17,19,43\) the denominator result ensures an affine reduction.
Enumerate the square residues \(b\) for which \(K_b\) splits into
squares. In each unramified factor of the descent algebra compare
the component norm of \(b-\theta\) with that of the representative,
allowing a common scalar as in \S\ref{arith:70}. A factor where \(b-\theta\)
vanishes is omitted from the test. The successive numbers of classes
are 6, 5, 4 and 3, and the last three are precisely the stated ones.
All three representatives are collisions by their explicit chain
vertices.
\cert{check_type8_product.m}

These classes give three explicit elliptic covers over the quartic
field \(K=\mathbf Q(\rho)\), \(G(\rho)=0\). Write
\(q(T)=G(T)/(1280(T-\rho))\). In the two-cover test the component at
zero is \(b\); since both \(b\) and \(b_0\) are nonzero squares, the
common rational scalar is square. Hence \((b-\rho)/(b_0-\rho)\)
is square over \(K\). As \(G(b)\) and \(G(b_0)\) are also squares,
the corresponding elliptic equation is \(W^2=q(b_0)q(b)\).
The next step uses the full singleton polynomial
and a second elliptic cover.

\subsubsection{A singular fiber and the reciprocal covers}\label{arith:73}

The class represented by \(b_0=1/4\) contains no
valid chain at any rational height. For the three classes of \S\ref{arith:72},
in the order \(b_0=1/4,9/4,25/4\), the quartic-field elliptic curves
there have 2-Selmer dimensions \(5,4,5\). Each class also gives an
elliptic curve with 2-Selmer dimension three. These computations are
unconditional.

\emph{Proof.} Scale the center and singleton squares by four. Write
\(B=4b\), \(B_0=4b_0\) and
\[P_B(X)=(X-B)((X-B-4)^2-16B)((X-B-16)^2-64B),\qquad
 F(B,X)=\frac{P_B(X)-P_B(1)}{X-1}.\]
All four roots of this monic quartic must be rational squares. Its
constant term is \(5f(B)\), where
\(f(T)=T^4-18T^3+114T^2-97T+405\). For a root \(r\) of \(f\),
\S\ref{arith:72} gives \((B-r)/(B_0-r)\in K^{\times2}\) in
\(K=\mathbf Q(r)\).

\textbf{The singular fiber.} \S\ref{arith:72} gives \(B\in\mathbf Z_{23}\),
so all four roots are integral there. Enumerating the twelve square
residues modulo 23, including zero, gives precisely \(B=1,2,9\)
for which all four roots are squares, counted with multiplicity.
Since \(-4\) is a simple root of \(f\bmod23\), the degree-one
completion there tests the squareclass by \((B+4)(B_0+4)\).
For \(B_0=1\) only \(B\equiv1\) survives. The special fiber factors as
\[F(1,X)=(X-1)(X-9)^2(X-25).\]
In particular some root satisfies \(X\equiv9\pmod{23}\).
Put \(U=B-1, V=X-9\). The terms of degree at most two of
\(F(1+U,9+V)\) are \(-64(2V^2-15U^2)\).
Since 30 is nonsquare modulo 23, this quadratic form has no
nontrivial zero over the residue field. If \((U,V)\ne(0,0)\),
put \(m=\min(v_{23}(U),v_{23}(V))\ge1\). Its quadratic term has
valuation exactly \(2m\), whereas all higher terms have valuation
at least \(3m\). Thus \(B=1\) and \(X=9\), and the repeated
singleton violates distinctness.
\cert{type8_descent/node.m}

\textbf{Two elliptic curves per class.} Write \(q(T)=f(T)/(T-r)\) and
\[g(T)=T^3q(1/T)/q(0).\]
Both cubics are monic and irreducible over \(K\). The direct
cover has equation \(W^2=q(B_0)q(B)\). Since \(B\) and \(B_0\)
are nonzero squares, the factor \((B B_0)^{-3}/q(0)^2\) is a
square as well. Thus the reciprocal cover is
\(W^2=g(1/B_0)g(1/B)\). For either monic cubic
\(h(T)=T^3+c_2T^2+c_1T+c_0\), writing \(k=h(t_0)\) gives
the Weierstrass model \(y^2=x^3+kc_2x^2+k^2c_1x+k^3c_0\),
with \(x=kt\). Irreducibility gives trivial rational two-torsion.

The cubic root field over \(K\) is the same for all six curves:
scaling a root by \(k\) or taking its reciprocal preserves the
field. It is the degree-twelve field \(M=\mathbf Q(r,r')\)
of two distinct ordered roots of \(f\). The certificate verifies
an absolute model, signature \((0,6)\), and discriminant
\(5528109^5\). Apply the explicit formula of \S\ref{sec:marked} with support
\(2^{22}\), generating bound \(2^{17}\) and no real places.
The rational interval bound gives
\(A_{\rm arch}+S-\log|D_M|>862/1000\).
Consequently the primes of norm at most \(2^{17}\) generate the
class group. PARI verifies their reduction to a factor base of
norms at most 1000. With this certified bound for \(M\) and the
Minkowski bound for other fields, the class group is cyclic of
order three and two-descent gives the dimensions asserted above.

\cert{type13_prym_class/prime_counts.cpp},
\cert{type13_prym_class/check_bound.py},
\cert{type8_descent/factor_base.gp},
\cert{type8_descent/covers.m}

\subsubsection{Global sign constraints and transverse logarithms}\label{arith:74}

No valid distinct-entry type-8 chain exists at any
rational height. In the classes represented by \(b_0=9/4,25/4\), the
reciprocal elliptic curves of \S\ref{arith:73} both have rank three.

\emph{Proof.} Continue with \(B=4b\), \(t=1/B\), \(t_0=1/B_0\),
\(k=g(t_0)\) and the monic reciprocal cubic \(g\) of \S\ref{arith:73}.
Write \(E\) for the model with abscissa \(x=kt\), and put
\(P_0=(kt_0,k^2)\). Use
\[K=\mathbf Q(h),\quad h^4-h^3+28h^2-54h+231=0,
 \qquad r=(-2h^3+h^2-33h+129)/15 .\]
For each branch the certificate supplies three points \(P_i\) of
\(E(K)\) by giving their parameters \(t_i\) and extracting a square
root of \(k^3g(t_i)\) in \(K\). One is \(P_0\); for \(B_0=9\) another
has rational parameter \(-59/64\). The logarithm calculation below
proves independence, so the upper bound three is attained in both cases.
\cert{type8_descent/logarithms.m}

\textbf{Signs for \(B_0=9\).} At each of \(p=17,131,197,269\), enumerate
square residues \(B\) for which \(F(B,X)\) splits into four squares.
For every monic irreducible factor \(d\) of \(f\bmod p\), require
\(d(B)d(9)\) to be square, allowing zero. This follows by taking
the residue-field norm of \((B-r)/(9-r)\) at the corresponding
unramified place when both entries are units. \S\ref{arith:72} excludes
poles at these primes. The sole surviving residue is \(B=9\) in each
case. All four primes give good reduction of \(E\). At each place
above them a putative global point therefore reduces to \(P_0\) or
\(-P_0\).

At 131 the residue degrees are \(1,1,2\). To determine the relative
signs, take the subgroup generated by the three independent points.
For \(\ell=2,3\) separately the certificate saturates this subgroup at
\(\ell\): for each line of coefficient vectors modulo \(\ell\), it
tests exact division of the indicated sum in \(E(K)\). When divisible,
replace the generator whose coefficient is one by the divided point
and repeat. This enlarges the subgroup by index \(\ell\). The final
tests have no further division and \(E(K)[\ell]=0\). Since the rank
is three, the resulting points represent all global points modulo
\(mE(K)\), with \(m=16\) or \(9\), respectively.

At each finite place with residue group of order \(N\), multiply by
\(N/\gcd(N,m)\) to annihilate \(m\) times the residue group. Enumerate
coefficient triples modulo \(m\) and require agreement with a signed
reduction of \(P_0\) at every place. For \(m=16\), the primes
\(17,197,131\) leave four triples, and the signs at the two linear
factors of 131 agree. For \(m=9\), the primes \(269,131\) leave
36 triples and tie the quadratic factor to one linear factor.
The only simultaneous patterns are \((+,+,+)\) and \((-,-,-)\).
Negating the global point if necessary makes all signs positive.
\cert{type8_descent/signs.m}

\textbf{Transversality.} Use \(p=131\) for \(B_0=9\) and \(p=2\) for
\(B_0=25\). In the second case there is just one place. \S\ref{arith:72}
makes \(B\) integral at 2. An even \(B\) would give
\(F(B,X)\equiv X^4+X^3+X^2+X+1\pmod2\), which has no root.
Thus \(B\) is an odd square and \(t\equiv t_0\pmod8\).

Let \(V=\prod_{v\mid p}K_v\), of dimension four over \(\mathbf Q_p\),
and take local elliptic logarithms for \(dx/(2y)\). The certificate
expresses these in integral power bases of the unramified completions.
For each \(P_i\) it takes a common integer multiple across the places,
computes the formal parameter \(z=-x/y\), and divides the concatenated
row by the same power of \(p\) at all places to make it primitive.
Every formal parameter used has valuation at least eight. For the
integral Weierstrass model the coefficient of \(z^j\) in the formal
logarithm lies in \(j^{-1}\mathcal O_{K_v}\), by integration of the
invariant differential. Hence replacing the logarithm by
\(z\) preserves each normalized coordinate modulo \(p^5\).

Let \(D\) be the resulting three-row matrix of true normalized logarithms
and let \(a\) be the coordinate vector of \(1/k\) in \(V\). For 131 the
minimum valuation of its maximal minors is zero, and for 2 it is one.
Adjoining \(a\) gives a determinant of exactly the same valuation in
both cases. These statements are determined by the calculation modulo
\(p^5\). In particular the rows are independent. Their signed maximal
minors, divided by the common power of \(p\), define an integral linear
functional \(\lambda\) annihilating all three rows, with
\(\lambda(1/k)\) a unit. As the three points span \(E(K)\otimes\mathbf Q\),
the functional kills the logarithm of every global point.

\textbf{The full parameter disks.} Write \(t=t_0+d\xi\), with
\(d=131\) or \(8\) and \(\xi\in\mathbf Z_p\). At 131 the sign result
puts the lift on the branch through \(P_0\) at every place. At 2 choose
the sign giving \(y/k^2\in1+4\mathcal O_{K_v}\).
On these branches \(g(t_0+d\xi)/k\) has constant term one and its
square root takes value one at zero. For \(t\ne t_0\) the chord used
in subtracting \(P_0\) has slope \((y+k^2)/(k(t-t_0))\). Its valuation
is \(-v_p(t-t_0)\) at 131 and \(1-v_2(t-t_0)\) at two. The difference
therefore belongs to the formal group throughout the disk. Pulling
back the differential, and dividing out the scalar \(d/2\), its
logarithm has the same zeros under \(\lambda\) as
\[\lambda\!\left(\frac1k\int_0^\xi
     \left(\frac{g(t_0+dX)}{k}\right)^{-1/2}dX\right).\]
The linear coefficient is a unit. At 131 the coefficient of \(X^i\)
in the integrand has valuation at least \(i\), since the binomial
coefficients of exponent \(-1/2\) are 131-integral. At two, for
\(i\ge1\), the valuation is at least \(i+1\): a term with \(j\)
nonconstant factors gains \(3i\) from \(d^i\), and
\(v_2\binom{-1/2}{j}=s_2(j)-2j\ge1-2i\), where \(s_2\) counts
binary ones. After division by \(i+1\) the bounds are still positive
and tend to infinity. Thus every nonlinear coefficient belongs to
\(p\mathbf Z_p\); for \(\xi\ne0\) the linear term has strictly smallest
valuation. Global points require \(\xi=0\).

The only parameters compatible with valid splitting in the two classes
are therefore \(B=9\) and \(25\), respectively. At both values
\(P_B(1)=0\), so the A and B vertex squares coincide. The class
\(B_0=1\) was excluded in \S\ref{arith:73}.
